\begin{derivation}[平面界面的混合表象与单位长度损失]
先从固定平行动量 $\boldsymbol\kappa_\parallel=(\kappa_x,\kappa_y)$ 的 $z$ 向标量 Green 核开始。介质 1 中它满足
\begin{equation*}
 (\partial_z^2+\kappa_{z1}^2)\mathscr G_z(z,z')
 =-\delta(z-z').
\end{equation*}
在 $z\neq z'$ 的区域，齐次解由推迟边界条件 $\operatorname{Im}\kappa_{z1}\ge0$ 选成
\begin{equation*}
 \mathscr G_z(z,z')
 =A\,\ee^{\ii\kappa_{z1}|z-z'|}.
\end{equation*}
为了确定 $A$，在 $z'\pm\epsilon$ 间积分微分方程：
\begin{align*}
 \int_{z'-\epsilon}^{z'+\epsilon}\dd z\,
 \partial_z^2\mathscr G_z
 +\kappa_{z1}^2
 \int_{z'-\epsilon}^{z'+\epsilon}\dd z\,\mathscr G_z
 &=-1.
\end{align*}
当 $\epsilon\to0$，第二项按 $O(\epsilon)$ 消失，所以
\begin{equation*}
 \partial_z\mathscr G_z\big|_{z'+0}
 -\partial_z\mathscr G_z\big|_{z'-0}=-1.
\end{equation*}
由
\begin{align*}
 \partial_z \ee^{\ii\kappa|z-z'|}
 &=\ii\kappa\,\operatorname{sgn}(z-z')
 \ee^{\ii\kappa|z-z'|}
\end{align*}
可得跳跃
\begin{equation*}
 2\ii\kappa_{z1}A=-1,
 \qquad
 A=\frac{\ii}{2\kappa_{z1}}.
\end{equation*}
这就是 Weyl 表示中每个平行动量通道的标量传播因子。

界面反射路径从源点 $z'$ 到界面再回到观察点 $z$，因此自由空间相位 $\ee^{\ii\kappa_{z1}|z-z'|}$ 变成
\begin{equation*}
 \ee^{\ii\kappa_{z1}(z+z')}.
\end{equation*}
对 TE、TM 两个物理偏振分别乘 Fresnel 反射系数 $r_s,r_p$，并插入对应的偏振二张量，就得到正文\eqnrefs{eq:planar_scattering_green_dyadic}。电子轨迹投影随后决定同步条件与有限长度线宽。

电子轨迹为
\begin{equation*}
 \mathbf r_e(t)=(vt,0,b),
 \qquad -\frac T2\le t\le\frac T2.
\end{equation*}
Weyl 相位含
\begin{equation*}
 \ee^{\ii\boldsymbol\kappa_\parallel\cdot
 [\mathbf R_e(t)-\mathbf R_e(t')]}
 =\ee^{\ii\kappa_xv(t-t')}.
\end{equation*}
频域 Green 泛函还含 $\ee^{\ii\omega(t'-t)}$，因此双时间积分的全部纵向相位组合成
\begin{equation*}
 \ee^{-\ii\Delta(t-t')},
 \qquad
 \Delta\equiv\omega-\kappa_xv.
\end{equation*}
双积分可先因子化：
\begin{align*}
 I_T(\Delta)
 &\equiv
 \int_{-T/2}^{T/2}\dd t\,
 \ee^{-\ii\Delta t}\\
 &=\left.
 \frac{\ee^{-\ii\Delta t}}{-\ii\Delta}
 \right|_{-T/2}^{T/2}\\
 &=\frac{2\sin(\Delta T/2)}{\Delta}\\
 &=T\,\sinc\!\left(\frac{\Delta T}{2}\right).
\end{align*}
于是
\begin{align*}
 &\int_{-T/2}^{T/2}\dd t
 \int_{-T/2}^{T/2}\dd t'\,
 \ee^{-\ii\Delta(t-t')}\\
 &\qquad=I_T(\Delta)I_T^*(\Delta)\\
 &\qquad=T^2\sinc^2\!\left(\frac{\Delta T}{2}\right).
\end{align*}
因此有限轨迹的同步线形不是 $\delta$ 函数，而是正文\eqnrefs{eq:finite-trajectory-sync-linewidth-dp83} 的 $\sinc^2$ 主峰。

第一零点来自
\begin{equation*}
 \sin\frac{\Delta T}{2}=0,
 \qquad
 \frac{|\Delta|T}{2}=\pi,
\end{equation*}
所以从中心到第一零点的角频率间隔是
\begin{equation*}
 \Delta\omega_{\rm zero}=\frac{2\pi}{T}.
\end{equation*}
半高点令
\begin{equation*}
 \sinc^2x=\frac12,
 \qquad x>0.
\end{equation*}
即求方程
\begin{equation*}
 \sin x=\frac{x}{\sqrt2}.
\end{equation*}
主峰内唯一正根为
\begin{equation*}
 x_{1/2}=1.391557\ldots.
\end{equation*}
左右半高点满足
\begin{equation*}
 \Delta_\pm=\pm\frac{2x_{1/2}}{T},
\end{equation*}
所以全高半宽
\begin{align*}
 \Delta\omega_{\rm FWHM}
 &=\Delta_+-\Delta_-\\
 &=\frac{4x_{1/2}}{T}\\
 &=\frac{5.56623\ldots}{T},
\end{align*}
即正文\eqnrefs{eq:finite-trajectory-sync-linewidth-dp83}。

再严格验证长时间极限。对任意在 $\Delta=0$ 附近连续且可积的测试函数 $f(\Delta)$，考虑
\begin{align*}
 J_T
 &\equiv\int_{-\infty}^{\infty}\dd\Delta\,
 f(\Delta)
 T\sinc^2\!\left(\frac{\Delta T}{2}\right).
\end{align*}
令 $x=\Delta T/2$，则 $\dd\Delta=2\dd x/T$：
\begin{align*}
 J_T
 &=2\int_{-\infty}^{\infty}\dd x\,
 f\!\left(\frac{2x}{T}\right)
 \sinc^2x.
\end{align*}
当 $T\to\infty$，对固定 $x$ 有 $f(2x/T)\to f(0)$，且
\begin{equation*}
 \int_{-\infty}^{\infty}\dd x\,\sinc^2x=\pi.
\end{equation*}
因此
\begin{equation*}
 \lim_{T\to\infty}J_T
 =2\pi f(0),
\end{equation*}
也就是分布恒等式
\begin{equation*}
 T\sinc^2\!\left(\frac{\Delta T}{2}\right)
 \xrightarrow[T\to\infty]{\mathcal D'}
 2\pi\delta(\Delta).
\end{equation*}

现在回到 Weyl 动量积分。长轨迹极限把
\begin{equation*}
 \delta(\omega-v\kappa_x)
\end{equation*}
插入 $\kappa_x$ 积分。由 $g(\kappa_x)=\omega-v\kappa_x$、$|g'|=v$，
\begin{equation*}
 \delta(\omega-v\kappa_x)
 =\frac1v
 \delta\!\left(\kappa_x-\frac\omega v\right).
\end{equation*}
因此
\begin{align*}
 \int_{-\infty}^{\infty}\frac{\dd\kappa_x}{2\pi}
 F(\kappa_x,\kappa_y)
 2\pi\delta(\omega-v\kappa_x)
 &=\frac1v
 F\!\left(\frac\omega v,\kappa_y\right).
\end{align*}
双时间积分提供的总因子是 $T$，而概率除以轨迹长度 $L=vT$；再乘上上式 Jacobian $1/v$ 与电流中的两个速度因子 $v^2$，有
\begin{equation*}
 \frac{T}{L}\times\frac1v\times v^2
 =1.
\end{equation*}
所以单位长度谱中不再残留额外的 $T$ 或 $v$ 前因子，$\kappa_x$ 被固定在同步值 $\omega/v$，只剩 $\kappa_y$ 积分。这正得到正文\eqnrefs{eq:planar_eels_mixed_green_per_length}。
\end{derivation}

\begin{derivation}[匀速电子的 Coulomb 场与跃迁矩阵元]
先取给定冲击参数下的直线匀速电子轨迹，并在这一计算层级忽略环境反作用对轨迹的改变。以靶中心为原点，轨迹写成 $(\mathbf b,z)$；相应实验室系 Liénard--Wiechert 电场为

\begin{align*}
\mathbf E_{e,\perp}(\mathbf b,z)
&=-\frac{e\gamma_0}{4\pi\varepsilon_0}
\frac{\mathbf b}{(b^2+\gamma_0^2z^2)^{3/2}},\\
E_{e,z}(\mathbf b,z)
&=-\frac{e\gamma_0}{4\pi\varepsilon_0}
\frac{z}{(b^2+\gamma_0^2z^2)^{3/2}}.
\end{align*}
电偶极相互作用为 $V(z)=-\mathbf d\cdot\mathbf E_e(z)$。盒平面波矩阵元含 Fourier 因子
\[
\frac1L\int_{-\infty}^{\infty}\dd z\,
\ee^{\ii\Delta\kappa z}V(z).
\]
令 $x=\gamma_0z$，则 $\dd z=\dd x/\gamma_0$，指数变成
\[
\ee^{\ii(\Delta\kappa/\gamma_0)x}.
\]
先以 $\kappa_x$ 表示一维傅里叶波数参数，单位 ${\rm m^{-1}}$。从基本积分
\[
\int_{-\infty}^{\infty}\frac{\ee^{\ii\kappa_x x}}{\sqrt{x^2+b^2}}\dd x
=2K_0(|\kappa_x|b)
\]
出发。对 $b$ 求导：
\begin{align*}
\frac{\partial}{\partial b}
\frac1{\sqrt{x^2+b^2}}
&=-\frac{b}{(x^2+b^2)^{3/2}},\\
\frac{\partial}{\partial b}2K_0(|\kappa_x|b)
&=-2|\kappa_x|K_1(|\kappa_x|b),
\end{align*}
所以
\[
\int\dd x\,
\frac{b\ee^{\ii\kappa_x x}}{(x^2+b^2)^{3/2}}
=2|\kappa_x|K_1(|\kappa_x|b).
\]
纵向积分利用
\[
\frac{x}{(x^2+b^2)^{3/2}}
=-\frac{\partial}{\partial x}
\frac1{\sqrt{x^2+b^2}}.
\]
分部积分，边界项在 $|x|\to\infty$ 消失：
\begin{align*}
\int\dd x\,
\frac{x\ee^{\ii\kappa_x x}}{(x^2+b^2)^{3/2}}
&=-\int\dd x\,\ee^{\ii\kappa_x x}
\frac{\partial}{\partial x}
\frac1{\sqrt{x^2+b^2}}\\
&=\ii\kappa_x\int\dd x\,
\frac{\ee^{\ii\kappa_x x}}{\sqrt{x^2+b^2}}\\
&=2\ii\kappa_xK_0(|\kappa_x|b).
\end{align*}
在积分恒等式中取参数 $|\Delta\kappa|/\gamma_0$，再分别乘回横向和纵向偶极分量以及盒归一化 $1/L$，即得\eqnrefs{eq:feberi_bessel_coulomb_matrix_element}。
\end{derivation}

\begin{derivation}[远场振幅与辐射能谱]
真空远场的 Fourier 分量满足

\begin{equation*}
 \widetilde{\mathbf H}_{\rm far}(\omega)
 =\frac1{\mu_0c}\hat{\mathbf n}\times
 \widetilde{\mathbf E}_{\rm far}(\omega),
 \qquad
 \hat{\mathbf n}\cdot\widetilde{\mathbf E}_{\rm far}=0.
\end{equation*}
时域总辐射能为
\begin{align*}
 \frac{\dd W_{\rm rad}}{\dd\Omega}
 &=r^2\int_{-\infty}^{\infty}\dd t\,
 \hat{\mathbf n}\cdot
 [\mathbf E_{\rm far}(t)\times\mathbf H_{\rm far}(t)]\\
 &=\frac{r^2}{\mu_0c}
 \int_{-\infty}^{\infty}\dd t\,|\mathbf E_{\rm far}(t)|^2.
\end{align*}
采用
\begin{equation*}
 \mathbf E(t)=\int\frac{\dd\omega}{2\pi}\,
 \widetilde{\mathbf E}(\omega)\ee^{-\ii\omega t},
\end{equation*}
Parseval 恒等式给
\begin{equation*}
 \int\dd t\,|\mathbf E(t)|^2
 =\int\frac{\dd\omega}{2\pi}|\widetilde{\mathbf E}(\omega)|^2.
\end{equation*}
实场满足 $\widetilde{\mathbf E}(-\omega)=\widetilde{\mathbf E}^{*}(\omega)$，因此正频部分为
\begin{equation*}
 \int\dd t\,|\mathbf E(t)|^2
 =\frac1\pi\int_0^\infty\dd\omega\,
 |\widetilde{\mathbf E}(\omega)|^2.
\end{equation*}
正文的远场源--场关系是
\begin{equation*}
 \widetilde{\mathbf E}_{\rm far}
 =\ii\mu_0\omega\frac{\ee^{\ii\kappa_0r}}{r}
 \mathbf A_\infty,
 \qquad \kappa_0=\omega/c.
\end{equation*}
代入 Poynting 谱后，$r^{-2}$ 与球面面积因子相消，得到
\begin{equation*}
 \frac{\dd^2 W_{\rm rad}}{\dd\omega\dd\Omega}
 =\frac{\mu_0\omega^2}{\pi c}|\mathbf A_\infty|^2,
\end{equation*}
即正文\eqnrefs{eq:far_field_radiated_energy_green_si}。正文\eqnrefs{eq:cl_photon_probability_far_green} 随后只是把这一能量谱除以单光子能量 $\hbar\omega$ 的定义转换，不需要另行引入动力学近似。
\end{derivation}

\begin{derivation}[自由空间电流的光子辐射概率]
真空标量 Helmholtz Green 函数
$g_0(\mathbf r-\mathbf r')=\ee^{\ii\kappa_0|\mathbf r-\mathbf r'|}/[4\pi|\mathbf r-\mathbf r'|]$ 的远场系数由严格极限定义：
\begin{equation*}
\lim_{r\to\infty}
4\pi r\,\ee^{-\ii\kappa_0r}
 g_0(r\hat{\mathbf n}-\mathbf r')
=\ee^{-\ii\kappa_0\hat{\mathbf n}\cdot\mathbf r'}.
\end{equation*}
电磁 Green 张量在同一极限只保留横向投影，因此
\[
\mathbf F_{\infty,0}
(\hat{\mathbf n},\mathbf r';\omega)
=\frac1{4\pi}
(\mathbf I-\hat{\mathbf n}\hat{\mathbf n})
\ee^{-\ii\kappa_0\hat{\mathbf n}\cdot\mathbf r'}.
\]
代入 $\mathbf A_\infty$ 定义：
\begin{align*}
\mathbf A_{\infty,0}
&=\int\dd^3 r'\,
\mathbf F_{\infty,0}\cdot\widetilde{\mathbf J}(\mathbf r',\omega)\\
&=\frac1{4\pi}
(\mathbf I-\hat{\mathbf n}\hat{\mathbf n})
\int\dd^3 r'\,\widetilde{\mathbf J}(\mathbf r',\omega)
\ee^{-\ii\kappa_0\hat{\mathbf n}\cdot\mathbf r'}\\
&=\frac1{4\pi}
(\mathbf I-\hat{\mathbf n}\hat{\mathbf n})
\widetilde{\mathbf J}(\kappa_0\hat{\mathbf n},\omega).
\end{align*}
因此
\[
|\mathbf A_{\infty,0}|^2
=\frac1{16\pi^2}
\left|(\mathbf I-\hat{\mathbf n}\hat{\mathbf n})
\widetilde{\mathbf J}\right|^2.
\]
代入\eqnrefs{eq:cl_photon_probability_far_green} 后，$1/\pi$ 与 $1/(16\pi^2)$ 合并成 $1/(16\pi^3)$，得到\eqnrefs{eq:free_space_current_radiation_probability}。
\end{derivation}

\begin{derivation}[能量损失、吸收与远场辐射的平衡]

先由外部电流对诱导场所做的负功求总损失。时域表达式为
\begin{equation*}
 W_{\rm loss}
 =-\int\dd t\dd^3 r\,
 \mathbf J(\mathbf r,t)\cdot\mathbf E(\mathbf r,t).
\end{equation*}
利用实场的正负频共轭关系，正频谱写成
\begin{equation*}
 \frac{\dd W_{\rm loss}}{\dd\omega}
 =-\frac1\pi\operatorname{Re}
 \int\dd^3 r\,
 \widetilde{\mathbf J}^{*}(\mathbf r,\omega)
 \cdot\widetilde{\mathbf E}(\mathbf r,\omega).
\end{equation*}
代入正文\eqnrefs{eq:eels_source_field_green_dp68}，
\begin{align*}
 \frac{\dd W_{\rm loss}}{\dd\omega}
 &=-\frac1\pi\operatorname{Re}
 \left[\ii\mu_0\omega
 \langle\widetilde J|G^R|\widetilde J\rangle\right]\\
 &=\frac{\mu_0\omega}{\pi}
 \langle\widetilde J|\operatorname{Im}G^R|\widetilde J\rangle,
\end{align*}
即正文\eqnrefs{eq:eels_energy_loss_green_si}。

再计算局域介质中的不可逆吸收。正文\eqnrefs{eq:material_polarization_current_dp68} 给出
\begin{equation*}
 \widetilde{\mathbf J}_{\rm mat}
 =-\ii\omega\varepsilon_0
 (\boldsymbol\varepsilon_r-\mathbf I)\widetilde{\mathbf E}.
\end{equation*}
取正频周期平均功，
\begin{align*}
 \frac{\dd W_{\rm abs}}{\dd\omega}
 &=\frac1\pi\int_{V_{\rm mat}}\dd^3 r\,
 \operatorname{Re}
 \left[\widetilde{\mathbf J}_{\rm mat}\cdot
 \widetilde{\mathbf E}^{*}\right]\\
 &=\frac{\omega\varepsilon_0}{\pi}
 \int_{V_{\rm mat}}\dd^3 r\,
 \widetilde{\mathbf E}^{*}\cdot
 \operatorname{Im}\boldsymbol\varepsilon_r
 \cdot\widetilde{\mathbf E},
\end{align*}
得到正文\eqnrefs{eq:material_absorbed_energy_spectrum_si}。

最后证明总损失如何严格分成远场辐射与材料吸收。取一个包围全部耗散材料、外边界位于无损真空的控制体 $V$。按全文 $\ee^{-\ii\omega t}$ 约定，
\begin{equation*}
 \nabla\times\mathbf E=\ii\omega\mathbf B,
 \qquad
 \nabla\times\mathbf H=\mathbf J_{\rm ext}-\ii\omega\mathbf D,
\end{equation*}
所以共轭方程为
\begin{equation*}
 \nabla\times\mathbf H^*=\mathbf J_{\rm ext}^*+\ii\omega\mathbf D^*.
\end{equation*}
代入恒等式
\begin{equation*}
 \nabla\cdot(\mathbf E\times\mathbf H^*)
 =\mathbf H^*\cdot(\nabla\times\mathbf E)
 -\mathbf E\cdot(\nabla\times\mathbf H^*)
\end{equation*}
得到复 Poynting 方程
\begin{equation*}
 \nabla\cdot(\mathbf E\times\mathbf H^*)
 =\ii\omega\mathbf H^*\cdot\mathbf B
 -\mathbf E\cdot\mathbf J_{\rm ext}^*
 -\ii\omega\mathbf E\cdot\mathbf D^*.
\end{equation*}
介质只含电损耗并取 $\mathbf B=\mu_0\mathbf H$，第一项为纯虚数，对实部没有贡献。再写
\begin{equation*}
 \mathbf D=\varepsilon_0\boldsymbol\varepsilon_r\mathbf E,
 \qquad
 X\equiv\mathbf E^*\cdot\boldsymbol\varepsilon_r\mathbf E,
\end{equation*}
则 $\mathbf E\cdot\mathbf D^*=\varepsilon_0X^*$，从而
\begin{equation*}
 \operatorname{Re}[-\ii\omega\mathbf E\cdot\mathbf D^*]
 =-\omega\varepsilon_0\operatorname{Im}X
 =-\omega\varepsilon_0
 \mathbf E^*\cdot\operatorname{Im}\boldsymbol\varepsilon_r\cdot\mathbf E.
\end{equation*}
因此
\begin{align*}
 \operatorname{Re}\nabla\cdot(\mathbf E\times\mathbf H^*)
 ={}&-\operatorname{Re}(\mathbf J_{\rm ext}^*\cdot\mathbf E)\\
 &-\omega\varepsilon_0
 \mathbf E^*\cdot\operatorname{Im}\boldsymbol\varepsilon_r\cdot\mathbf E.
\end{align*}
对 $V$ 积分并用 Gauss 定理，
\begin{align*}
 \operatorname{Re}\oint_{\partial V}\dd\mathbf S\cdot
 (\mathbf E\times\mathbf H^*)
 ={}&-\operatorname{Re}\int_V\dd^3r\,
 \mathbf J_{\rm ext}^*\cdot\mathbf E\\
 &-\omega\varepsilon_0\int_{V_{\rm mat}}\dd^3r\,
 \mathbf E^*\cdot\operatorname{Im}\boldsymbol\varepsilon_r\cdot\mathbf E.
\end{align*}
全文的正频 Fourier 归一化给每一项共同的 $1/\pi$。于是右边第一项正是前面得到的 $\dd W_{\rm loss}/\dd\omega$，第二项是 $-\dd W_{\rm abs}/\dd\omega$；把 $\partial V$ 推到无损真空远场后，左边就是 $\dd W_{\rm rad}/\dd\omega$。因此逐频有
\begin{equation*}
 \frac{\dd W_{\rm loss}}{\dd\omega}
 =\frac{\dd W_{\rm rad}}{\dd\omega}
 +\frac{\dd W_{\rm abs}}{\dd\omega},
\end{equation*}
即正文\eqnrefs{eq:eels_radiative_absorptive_energy_balance}。这里的分解来自同一个 Maxwell 能量守恒恒等式，而不是把 EELS 与 CL 当成两套独立动力学。
\end{derivation}

\begin{derivation}[平面反射 Green 张量的纵向分量]
真空侧平面波的单位 TE 偏振定义为

\[
\hat{\mathbf s}
=\frac{\hat{\mathbf z}\times\mathbf \kappa_\parallel}{\kappa_\parallel}.
\]
若
\[
\mathbf \kappa_\parallel=\kappa_x\hat{\mathbf x}+\kappa_y\hat{\mathbf y},
\]
则
\[
\hat{\mathbf s}
=\frac{-\kappa_y\hat{\mathbf x}+\kappa_x\hat{\mathbf y}}{\kappa_\parallel},
\]
因此
\[
|\hat{\mathbf x}\cdot\hat{\mathbf s}|^2
=\frac{\kappa_y^2}{\kappa_\parallel^2}.
\]

真空 TM 偏振取
\[
\hat{\mathbf p}_1^\pm
=\frac{\kappa_\parallel\hat{\mathbf z}
\mp \kappa_{z1}\hat{\boldsymbol\kappa}_\parallel}{\kappa_0},
\qquad
\hat{\boldsymbol\kappa}_\parallel
=\frac{\mathbf \kappa_\parallel}{\kappa_\parallel}.
\]
所以
\begin{align*}
\hat{\mathbf x}\cdot\hat{\mathbf p}_1^+
&=-\frac{\kappa_{z1}}{\kappa_0}\frac{\kappa_x}{\kappa_\parallel},\\
\hat{\mathbf p}_1^-\cdot\hat{\mathbf x}
&=+\frac{\kappa_{z1}}{\kappa_0}\frac{\kappa_x}{\kappa_\parallel},
\end{align*}
从而
\[
(\hat{\mathbf x}\cdot\hat{\mathbf p}_1^+)
(\hat{\mathbf p}_1^-\cdot\hat{\mathbf x})
=-\frac{\kappa_{z1}^2}{\kappa_0^2}
\frac{\kappa_x^2}{\kappa_\parallel^2}.
\]
亚光速同步条件给 $\kappa_\parallel>\kappa_0$，故
\[
\kappa_{z1}=\ii\alpha_1,
\qquad
-\kappa_{z1}^2=\alpha_1^2.
\]
于是 TM 投影化为
\[
\frac{\alpha_1^2}{\kappa_0^2}
\frac{\kappa_x^2}{\kappa_\parallel^2}.
\]
Weyl 散射 Green 张量的公共传播因子为
\[
\frac{\ii}{2\kappa_{z1}}
\ee^{\ii \kappa_{z1}(z+z')}.
\]
取源点和观察点都在电子轨迹 $z=z'=b$：
\begin{align*}
\frac{\ii}{2\kappa_{z1}}\ee^{2\ii \kappa_{z1}b}
&=\frac{\ii}{2\ii\alpha_1}\ee^{-2\alpha_1 b}\\
&=\frac{\ee^{-2\alpha_1 b}}{2\alpha_1}.
\end{align*}
TE 投影乘该因子给
\[
\ee^{-2\alpha_1 b}\frac{\kappa_y^2}{2\alpha_1 \kappa_\parallel^2}r_s,
\]
TM 投影乘该因子给
\[
\ee^{-2\alpha_1 b}
\frac{\alpha_1 \kappa_x^2}{2\kappa_0^2\kappa_\parallel^2}r_p.
\]
两项相加得到\eqnrefs{eq:planar_aloof_gxx_explicit}。
\end{derivation}

\begin{derivation}[平行掠射电子的 Fresnel 损失积分]
三维 Green 张量满足

\[
[\nabla\times\nabla\times-\kappa_0^2\varepsilon]\mathbf G^R
=\mathbf I\delta^{(3)}(\mathbf r-\mathbf r'),
\]
因此 $[\mathbf G^R]={\rm m^{-1}}$。Weyl 表示定义为
\[
\mathbf G^R
=\int\frac{\dd^2 \kappa_\parallel}{(2\pi)^2}
\ee^{\ii\mathbf \kappa_\parallel\cdot(\mathbf R-\mathbf R')}
\mathbf G_{{\rm sc},W}^R.
\]
因为 $[\dd^2 \kappa_\parallel]={\rm m^{-2}}$，必须有
\[
[\mathbf G_{{\rm sc},W}^R]={\rm m}.
\]
因此 $\dd \kappa_y\,G_{{\rm sc},W;xx}^R$ 无量纲。

前因子单位为
\begin{align*}
\left[\frac{\mu_0e^2}{\hbar}\right]
&=\frac{({\rm kg\,m\,C^{-2}})({\rm C^2})}
{{\rm J\,s}}\\
&=\frac{{\rm kg\,m}}{{\rm kg\,m^2\,s^{-1}}}\\
&={\rm s\,m^{-1}}.
\end{align*}
所以
\[
\left[\frac1L\frac{\dd P}{\dd\omega}\right]
={\rm s\,m^{-1}}.
\]
若换变量 $E_{\rm loss}=\hbar\omega$，
\[
\frac{\dd\omega}{\dd E_{\rm loss}}=\frac1\hbar,
\]
因此
\[
\frac1L\frac{\dd P}{\dd E_{\rm loss}}
=\frac1\hbar\frac1L\frac{\dd P}{\dd\omega}.
\]
\end{derivation}

\begin{derivation}[Drude 表面模与电子相位同步]
无损 Drude 介电函数为

\[
\varepsilon_2(\omega)=1-\frac{\omega_p^2}{\omega^2}.
\]
真空--介质界面 TM 表面模色散满足
\[
\frac{\kappa_{\rm sp}}{\kappa_0}
=\sqrt{\frac{\varepsilon_2}{1+\varepsilon_2}}.
\]
在 $\kappa_y=0$ 时，电子源的同步条件是
\[
\kappa_{\rm sp}=\frac{\omega}{v}
=\frac{\kappa_0}{\beta_{\rm rel}}.
\]
平方并消去 $\kappa_0$：
\[
\frac1{\beta_{\rm rel}^2}
=\frac{\varepsilon_2}{1+\varepsilon_2}.
\]
交叉相乘：
\begin{align*}
1+\varepsilon_2
&=\beta_{\rm rel}^2\varepsilon_2,\\
1
&=(\beta_{\rm rel}^2-1)\varepsilon_2,\\
\varepsilon_2
&=-\frac1{1-\beta_{\rm rel}^2}.
\end{align*}
再代入 Drude 形式：
\begin{align*}
1-\frac{\omega_p^2}{\omega^2}
&=-\frac1{1-\beta_{\rm rel}^2},\\
\frac{\omega_p^2}{\omega^2}
&=1+\frac1{1-\beta_{\rm rel}^2}\\
&=\frac{2-\beta_{\rm rel}^2}{1-\beta_{\rm rel}^2}.
\end{align*}
所以
\[
\omega^2
=\omega_p^2\frac{1-\beta_{\rm rel}^2}{2-\beta_{\rm rel}^2}.
\]
取正频率根得到\eqnrefs{eq:drude_surface_pole_electron_sync_frequency}。
\end{derivation}
